%%%%%%%%%%%%%%%%%%%%%%%%%%%%%%%%%
% PACKAGE IMPORTS
%%%%%%%%%%%%%%%%%%%%%%%%%%%%%%%%%


\usepackage[tmargin=2cm,rmargin=1in,lmargin=1in,margin=0.85in,bmargin=2cm,footskip=.2in]{geometry}
\usepackage[utf8]{inputenc}
\usepackage[T1]{fontenc}
\usepackage{amsmath,amsfonts,amsthm,amssymb,mathtools}
\usepackage{bookmark}
\usepackage[shortlabels]{enumitem}
\usepackage{hyperref,theoremref}
\usepackage{xcolor}
\hypersetup{
    pdftitle={Probability III Class Notes},
    colorlinks=true, linkcolor=blue!90,
    bookmarksnumbered=true,
    bookmarksopen=true
}

\usepackage{tikz}
\usetikzlibrary{automata}
\usetikzlibrary{positioning, arrows}

%\usepackage{import}
%\usepackage{xifthen}
%\usepackage{pdfpages}
%\usepackage{transparent}


%%%%%%%%%%%%%%%%%%%%%%%%%%%%%%
% SELF MADE COLORS
%%%%%%%%%%%%%%%%%%%%%%%%%%%%%%


\usepackage{anyfontsize}
% --- amsthm Configuration ---

% Style 1: The "Plain" style (Bold header, Italic body)
% Best for Theorems, Lemmas, and Corollaries.
\theoremstyle{plain}
\newtheorem{thm}{Theorem}[section]
\newtheorem{prop}[thm]{Proposition}

% Style 2: The "Definition" style (Bold header, Roman body)
% Best for Definitions and Examples.
\theoremstyle{definition}
\newtheorem{defn}{Definition}[section]
\newtheorem{exmp}{Example}[section]

% Style 3: The "Remark" style (Italic header, Roman body)
% Best for Remarks and Notes.
\theoremstyle{remark}
\newtheorem*{remark}{Remark}
\newtheorem{question}{Question}
